\section{Introduction}\label{sec:introduction}

MHC class I molecules present 8--11-mer peptides to CD8+ T cells, and accurate prediction of peptide--MHC-I binding underpins epitope discovery for vaccine design and cancer immunotherapy \citep{Falk91,Lundegaard11}. For HLA-A*02:01 --- the most prevalent class I allele globally --- binding is dominated by two primary anchor positions: position 2 (p2, hydrophobic: Leu/Met/Ile/Val) and position 9 (p9, C-terminal, Val/Leu preferred) \citep{Ruppert93,Rammensee95}. Computational predictors have evolved from position-specific scoring matrices to neural networks and pan-specific tools such as NetMHCpan \citep{Reynisson20}, MHCflurry \citep{ODonnell20}, BigMHC \citep{Albert23}, and MUNIS \citep{Chen25}, increasingly using protein language model embeddings \citep{Rives21}.

Several gaps remain. Head-to-head architecture comparisons under fully controlled conditions (identical training data, encoding, evaluation) are uncommon. The impact of training-label quality on downstream performance is seldom quantified, despite evidence that $>$55\% of IEDB entries are computationally rather than experimentally labelled \citep{Preibisch26}. And the translation of binding prediction to real-world epitope discovery is rarely demonstrated in a unified, reproducible pipeline.

We address these gaps through a systematic benchmark designed to isolate the effects of architecture choice, label quality, and encoding. Our contributions are: (i) quantification of label-quality effects across three labelling strategies, showing label provenance accounts for more performance variation than architecture choice; (ii) a controlled six-architecture comparison on a common dataset; (iii) evaluation of ESM-2 embeddings against BLOSUM62 with a mean-pooled ablation; (iv) external validation against 49 IEDB epitopes; (v) translational application through protein scanning and cancer mutation analysis; (vi) a docking-based structural hypothesis for non-canonical anchoring; and (vii) a multi-allele extension testing whether the recipe generalizes across common HLA-I alleles and how hard non-canonical binders are for sequence- and structure-based predictors alike.
